\documentclass[10pt,a4paper]{amsart}
\usepackage[margin=19mm]{geometry}
\usepackage{amsmath,amssymb,mathtools}
\usepackage[scr=rsfs,cal=euler]{mathalfa}
\usepackage{xfrac}
\usepackage[unicode,hidelinks,bookmarks=false]{hyperref}
\newcommand{\ie}{i.e.\ }
\newcommand{\cf}{cf.\ }
\newcommand{\resp}{respectively\ }
\newcommand{\Subsection}{\subsection}
\providecommand{\nobreakdash}{\nobreak\hbox{-}}
\setlength{\emergencystretch}{2em}
\multlinegap0pt
\usepackage{amssymb}
\usepackage{enumerate}
\usepackage[shortlabels]{enumitem}
\newtheorem{proposition}{Proposition}
\newcommand{\Ltwomap}{\mathfrak{s}}
\newcommand{\cofu}[1]{\mathop{\boxtimes}_{#1}}
\newcommand{\rar}{\rightarrow}
\newcommand{\tens}[1]{\mathop{*}\limits_{#1}}
\title{A three-functor formalism for commutative von Neumann algebras}
\author{Andre G. Henriques, Thomas A. Wasserman}
\date{Source: arXiv:2504.10131v2 (CC BY 4.0)}
\begin{document}
\maketitle

\noindent\emph{Source and licence.} A three-functor formalism for commutative von Neumann algebras. Version arXiv:2504.10131v2 (CC BY 4.0), distributed by arXiv under the Creative Commons Attribution 4.0 International licence, http://creativecommons.org/licenses/by/4.0/, as recorded in the arXiv licence metadata for this exact version and shown on the abstract page https://arxiv.org/abs/2504.10131v2. Full licence notice: \texttt{licenses/arxiv-2504.10131-CC-BY-4.0.txt}.\par\smallskip

\noindent\textit{Editorial scope.} Counted unit: the article's proposition prop:tensstandard (source0.tex lines 912--918). The article's complete original proof of it is delivered with it (source0.tex lines 920--979, one \texttt{proof} environment of 60 lines). No mathematics is written, repaired or re-derived: the statement and the proof are the article's own lines, in the article's own notation, and no retained line is substituted or re-flowed.  Retained uncounted as the source-local context the proof needs: lines 758--764; lines 883--885.  Nothing was omitted from any retained span, so the package carries no omission record.  No retained line contains a citation, so the wrapper carries no bibliography and no bibliography entry.  No retained line contains a figure environment, an image-inclusion command or drawing code, so no image is delivered and the package declares no asset dependency.  Editorial formatting only, in addition: an AMS \texttt{amsart} preamble that restates the article's own declarations and macros used by the retained lines, namely the environments \texttt{proposition} on separate counters, together with the standard packages those lines need.  The retained environments print in this document as Proposition 1 in order of appearance, whereas the article prints the same items with the numbering of its own full document.  The complete native source of the article as distributed in the arXiv e-print for this exact version is bundled unchanged as \texttt{sources/176499/source0.tex}.
\begin{equation}\label{e:innerprodonfusion}
    \langle
\phi_1\otimes n_1,\phi_2\otimes n_2
\rangle:=\langle (\phi_2^*\phi_1)
n_1,n_2
\rangle,
\end{equation} 

\begin{equation}\label{last eq}
\sqrt{\phi}^* a \sqrt{\phi} = \phi(a).
\end{equation}

\begin{proposition}\label{prop:tensstandard}
Let $C\rar A$ and $C\rar B$ be maps of commutative von Neumann algebras. Then the map $\Ltwomap\colon 	L^2(A)\cofu{C}L^2(B) \rar  L^2(A\ast_{C}B)$ defined by
	$$
	\Ltwomap\colon 	\sqrt{\phi} \cofu{\mu}\sqrt{\psi} \mapsto 	\sqrt{\mu (\phi \otimes \psi)}
	$$
	is an isometric isomorphism (i.e.~unitary).
\end{proposition}

\begin{proof}
Write $C$ as $C=\bigoplus C_i$ where each $C_i$ admits a faithful state, and consider the corresponding decompositions $A=\bigoplus A_i$ and $B=\bigoplus B_i$, with maps $C_i \rar A_i$ and $C_i \rar B_i$. We can further decompose $A_i=\bigoplus A_{ij}$ where each $A_{ij}$ admits a faithful conditional expectation $A_{ij} \rar C_i$, and similarly for $B_i$.
This reduces the statement of the proposition to the special case when $C$ admits faithful states, and $A$ and $B$ admit faithful conditional expectations to $C$.\footnote{When the von Neumann algebras $A$, $B$, and $C$ are representable on separable Hilbert spaces, then faithful conditional expectations and states always exist, and there is no need to consider direct sum decompositions.}

Let $\phi:A\to C$ and $\psi:B\to C$ be faithful conditional expectations, and let $\mu:C\to \mathbb C$ be a faithful state, so that
$\sqrt{\phi} \cofu{\mu}\sqrt{\psi}\in L^2(A)\cofu{C}L^2(B)$. Then $\mu \phi$ and $\mu \psi$ are faithful states on $A$ and $B$ respectively, and $\mu(\phi \otimes\psi)$ is a faithful state on $A \ast_{ C} B$.
We first define a map
\begin{align*}
     \Ltwomap_{\phi,\mu,\psi}\colon& L^2(A)\cofu{C}L^2(B) \rar  L^2(A\tens{C}B)\\
    & (a\sqrt{\phi})\cofu{\mu}(b\sqrt{\psi})\mapsto a\otimes b \sqrt{\mu (\phi \otimes \psi)}
\end{align*}
which a priori depends on $\phi$, $\psi$, and $\mu$.
We claim that $\Ltwomap_{\phi,\mu,\psi}$ is a unitary isomorphism, and that it agrees with $\Ltwomap$.

To see that $ \Ltwomap_{\phi,\mu,\psi}$ is isometric, observe that
$$
\big\|\Ltwomap_{\phi,\mu,\psi}\big((a\sqrt{\phi})\cofu{\mu}(b\sqrt{\psi})\big)\big\|^2=\mu (\phi(a^*a)\psi(b^*b)),
$$
and
\begin{align*}
\big\|(a\sqrt{\phi})\cofu{\mu}(b&\sqrt{\psi})\big\|^2 =
\big\langle
a\sqrt{\phi}\otimes b\sqrt{\mu\psi},a\sqrt{\phi}\otimes b\sqrt{\mu\psi}
\big\rangle
\\[-1mm]
&\quad\,\,\,\,\stackrel{\eqref{e:innerprodonfusion}}=
\big\langle
(\sqrt{\phi}^*a^*a\sqrt{\phi}) b\sqrt{\mu\psi}, b\sqrt{\mu\psi}
\big\rangle
\\&\stackrel{\eqref{last eq}}=
\big\langle
(\phi(a^*a) b\sqrt{\mu\psi}, b\sqrt{\mu\psi}
\big\rangle
=
\mu\psi(\phi(a^*a)b^*b)=\mu (\phi(a^*a)\psi(b^*b)),
\end{align*}
where the last equality uses the $C$-linearity of $\psi$. The map $\Ltwomap_{\phi,\mu,\psi}$ is unitary as it induces a bijection between the dense subsets $A\odot_C B$ on both sides.

To see that $ \Ltwomap_{\phi,\mu,\psi}$ agrees with $\Ltwomap$, it suffices to consider the values of both maps on vectors of the form $\sqrt{\phi'} \cofu{\mu'}\sqrt{\psi'}$ where $\sqrt{\phi'}=a\sqrt{\phi}$, $\sqrt{\psi'}=b\sqrt{\psi}$ and $\sqrt{\mu'}=c\sqrt{\mu}$, and $a$, $b$ and $c$ are positive elements of $A$, $B$ and $C$, respectively. We have
$$
\Ltwomap(\sqrt{\phi'} \cofu{\mu'}\sqrt{\psi'})=\sqrt{\mu'(\phi' \otimes \psi')},
$$
while $ \Ltwomap_{\phi,\mu,\psi}$ sends 
$$
\sqrt{\phi'} \cofu{\mu'}\sqrt{\psi'}=a\sqrt{\phi} \cofu{\mu}cb\sqrt{\psi}\quad\text{to}\quad  a\otimes cb \sqrt{\mu(\phi \otimes \psi)}.
$$
As both $\sqrt{\mu'(\phi' \otimes \psi')}$ and $a\otimes cb \sqrt{\mu(\phi \otimes \psi)}$ are positive elements of $L^2(A \ast_{C}B)$, to see that they are equal, it is sufficient to check that the corresponding states on $A \ast_{C}B$ agree, namely:
\[
\big(a\otimes cb \sqrt{\mu(\phi \otimes \psi)}\;\!\big)^2=\mu'(\phi' \otimes \psi').
\]
Indeed, for $x\otimes y \in A \ast_{C}B$, one checks:
    \begin{align*}
    \!\!&\!\!\big(a\otimes cb \sqrt{\mu(\phi \otimes \psi)}\;\!\big)^2(x\otimes y)=
        \mu(\phi \otimes \psi)\big( (x\otimes y)(a\otimes cb)^2\big) \\
    \!\!&\!\!=\mu(\phi(x a^2)\psi(y c^2b^2))
        =\mu(\phi'(x)\psi'(y)c^2)
         =\mu'(\phi'(x)\psi'(y))=
         \big(\mu'(\phi' \otimes \psi')\big)(x\otimes y).\qedhere
    \end{align*}
\end{proof}

\end{document}
